%%%%%
%
\section{Correction des deltaFF (vents mesur�s par le dispositif flux a\'erodynamique)}

\itemstyle{Op\'erateur~:} {\sf Cor\_deltaFF}
\index{Cor\_deltaFF}

%%%%%OBJET%%%%
\itemstyle{Objet~:}
Cette fonction corrige les incoh\'erences dans les valeurs de vitesses
de vent mesur\'es (niveau bas et niveau haut) � partir du flag cr�� par l'op�rateur Valid\_empir\_deltaFF (voir doc de cet op�rateur). 
Si le flag est �gal � 0, les mesures de vent haut et bas ont �t� d�tect�es coh�rentes entr'elles. Si le flag est n�gatif, ffhaut est trop "fort" par rapport � ffbas (ffbas trop "faible"...). Si le flag est positif, ffhaut est trop faible...
Le cas le plus fr�quemment rencontr� est une mesure de vent bas (ffbas) trop "faible" (donc flag<0) , souvent l'an�mom�tre est bloqu� � la valeur offset (flag=-1).\\
Selon la valeur du flag, ffbas ou ffhaut sont corrig�s par m�thode empirique ou par valeur interpol�e (cas de points incoh�rents "isol�s").

%%%%%ARGUMENTS%%%%
\bigskip
%%%%%ENTREES%%%%
\begin{tabarg}
\multicolumn{3}{l}{\titretab {Entr\'ees}}\\
\hline
$ffbas$ & Vecteur, {\tt float} & Force du vent mesur\'e au niveau bas\\
$ffhaut$ & Vecteur, {\tt float} & Force du vent mesur\'e au niveau haut\\
$deltaff$ & Vecteur, {\tt float} & Gradient de vent : haut-bas\\
$flag$ & Vecteur, {\tt float} & Flag de coh�rence des mesures de vent (haut/bas)\\
$offset_{bas}$ & Coeff., {\tt float} & Valeur min du vent niveau bas\\
$offset_{haut}$ & Coeff., {\tt float} & Valeur maxi du vent niveau haut\\
\end{tabarg}

\medskip
%%%%%SORTIES%%%%
\begin{tabarg}
\multicolumn{3}{l}{\titretab {Sorties}}\\
\hline
$ffbas\_Cdf$ & Vecteur, {\tt float} & Force du vent au niveau bas valid� ou corrig�\\
$ffhaut\_Cdf$ & Vecteur, {\tt float} & Force du vent au niveau haut valid� ou corrig�\\
$deltaff\_Cdf$ & Vecteur, {\tt float} & Gradient de vent : haut-bas valid� ou corrig�\\
$flag\_Cdf$ & Vecteur, {\tt float} & Flag de qualification (0=ok, 1=interpolation, 2= correction empirique)\\
\end{tabarg}


%%%%%FORMULE%%%%
\itemstyle{Formule}

%%Cor_deltaFF
Selon la valeur du flag :\\
1- si le flag est �gal � 0 (mesures ok) alors : \\
\hspace*{2cm} $flag\_Cdf$ = 0 et $ffbas\_Cdf$ = $ffbas$ etc...
\\
\\
2- si le flag est positif ($ffhaut$ trop faible par rapport � $ffbas$) alors : \\
\\
\indent Si le point est isol� (flag avant et apr�s �gaux � 0) , on recalcule $ffhaut\_Cdf$ � partir des valeurs interpol�s selon la formule suivante : \\
\hspace*{2cm} $coeff$ = (( $ffhaut_{i-1}$ / $ffbas_{i-1}$ ) + ( $ffhaut_{i+1}$/$ffbas_{i+1}$ )) / 2.  \\
\hspace*{2cm} $ffhaut\_Cdf$ = ( 1. / $coeff$ ) * $ffbas$  \\
\hspace*{2cm} $ffbas\_Cdf$ = $ffbas$  \\
\hspace*{2cm} $flag\_Cdf$ = 1 \\
\\
\indent Si le point n'est pas isol� (flag avant et/ou apr�s diff�rent de 0) , on recalcule $ffhaut\_Cdf$ par la m�thode empirique qui a servi � d�tecter les points de mesures incoh�rents entr'eux : \\
\hspace*{2cm} $ffhaut\_Cdf$ = ( 1.1647 * $ffbas$ ) + 0.1  \\
\hspace*{2cm} $ffbas\_Cdf$ = $ffbas$  \\
\hspace*{2cm} $flag\_Cdf$ = 2 \\
\\
\\
3- si le flag est n�gatif ($ffbas$ trop faible par rapport � $ffhaut$) alors : \\
\\
\indent Si le point est isol� (flag avant et apr�s �gaux � 0) , on recalcule $ffbas\_Cdf$ � partir des valeurs interpol�s selon la formule suivante : \\
\hspace*{2cm} $coeff$ = (( $ffhaut_{i-1}$ / $ffbas_{i-1}$ ) + ( $ffhaut_{i+1}$/$ffbas_{i+1}$ )) / 2.  \\
\hspace*{2cm} $ffbas\_Cdf$ = $coeff$ * $ffbas$  \\
\hspace*{2cm} $ffhaut\_Cdf$ = $ffhaut$  \\
\hspace*{2cm} $flag\_Cdf$ = 1 \\
\\
\indent Si le point n'est pas isol� (flag avant et/ou apr�s diff�rent de 0) , on recalcule $ffbas\_Cdf$ par la m�thode empirique qui a servi � d�tecter les points de mesures incoh�rents entr'eux : \\
\hspace*{2cm} $ffbas\_Cdf$ = ( 0.8586 * $ffhaut$ ) - 0.08586  \\
\hspace*{2cm} $ffhaut\_Cdf$ = $ffhaut$  \\
\hspace*{2cm} $flag\_Cdf$ = 2 \\
\\
\\
RAPPEL : Formules des droites empiriques d�limitant la zone de coh�rence utilis�es pour la "d�tection" des anomalies (op�rateur Valid\_empir\_deltaFF) :\\
\hspace*{1cm} Si $ffbas$ < 4 m/s :\\
\hspace*{2cm} $ffhaut$ doit �tre inf�rieur � $ffhaut_{maxi1}$ avec \\ 
$ffhaut_{maxi1}$ = 1.225 * $ffbas$ + 0.7\\
\\
\hspace*{1cm} Sinon :\\
\hspace*{2cm} $ffhaut$ doit �tre inf�rieur � $ffhaut_{maxi2}$ avec \\ 
$ffhaut_{maxi2}$ = 1.6 * $ffbas$ - 0.8\\
\\
\hspace*{1cm} Dans tous les cas (quelle que soit la valeur de $ffbas$):\\
\hspace*{2cm} $ffhaut$ doit �tre sup�rieur � $ffhaut_{min}$ avec \\ 
$ffhaut_{min}$ = 1.08 * $ffbas$ - 0.5\\


%%%%%COMMENTAIRE%%%%
\itemstyle{Commentaire~:}
Cette fonction ne permet que de corriger un probl�me de coh\'erence des
mesures (haut et bas) entre elles. Elle ne donne pas de      
diagnostic absolu de la validit\'e des mesures en elles-m�mes.

%%%%%AUTEUR%%%%
\itemstyle{Auteur~:}
CNRM/GMEI/TRAMM

%%%%%REFERENCE%%%%
\itemstyle{R\'ef\'erences~:}

